% =============================================================================
% VisAi — Analisi dei Requisiti Tecnici
% Progetto: Sistema Anti-Taccheggio basato su Computer Vision ed Edge AI
% Autore: Andrea Masiero
% Data: Luglio 2026
% =============================================================================

\documentclass[a4paper, 12pt]{article}

% --- Pacchetti ---
\usepackage[utf8]{inputenc}
\usepackage[T1]{fontenc}
\usepackage[italian]{babel}
\usepackage{geometry}
\usepackage{graphicx}
\usepackage{xcolor}
\usepackage{hyperref}
\usepackage{booktabs}
\usepackage{tabularx}
\usepackage{longtable}
\usepackage{array}
\usepackage{fancyhdr}
\usepackage{titlesec}

\usepackage{amsmath}
\usepackage{tikz}
\usetikzlibrary{calc}
\usetikzlibrary{shapes.geometric, arrows.meta, positioning, fit, backgrounds}

% --- Geometria pagina ---
\geometry{
    top=2.5cm,
    bottom=2.5cm,
    left=2.5cm,
    right=2.5cm
}

% --- Colori personalizzati ---
\definecolor{visaiblue}{HTML}{1A73E8}
\definecolor{visaidark}{HTML}{1B1F3B}
\definecolor{visaigray}{HTML}{5F6368}
\definecolor{visailight}{HTML}{E8F0FE}
\definecolor{visaigreen}{HTML}{34A853}
\definecolor{visaiorange}{HTML}{F9AB00}
\definecolor{visaired}{HTML}{EA4335}

% --- Hyperref ---
\hypersetup{
    colorlinks=true,
    linkcolor=visaiblue,
    urlcolor=visaiblue,
    citecolor=visaiblue,
    pdftitle={VisAi — Analisi dei Requisiti Tecnici},
    pdfauthor={Andrea Masiero}
}

% --- Header/Footer ---
\pagestyle{fancy}
\fancyhf{}
\fancyhead[L]{\footnotesize\textcolor{visaigray}{VisAi — Analisi dei Requisiti Tecnici}}
\fancyhead[R]{\footnotesize\textcolor{visaigray}{Luglio 2026}}
\fancyfoot[C]{\thepage}
\renewcommand{\headrulewidth}{0.4pt}

% --- Titoli sezione ---
\titleformat{\section}
    {\Large\bfseries\color{visaidark}}
    {\thesection}{1em}{}[\titlerule]
\titleformat{\subsection}
    {\large\bfseries\color{visaiblue}}
    {\thesubsection}{1em}{}
\titleformat{\subsubsection}
    {\normalsize\bfseries\color{visaidark}}
    {\thesubsubsection}{1em}{}

% --- Comandi personalizzati ---
\newcommand{\reqid}[1]{\texttt{\textcolor{visaiblue}{#1}}}
\newcommand{\priority}[1]{%
    \ifx#1H\textcolor{visaired}{\textbf{ALTA}}%
    \else\ifx#1M\textcolor{visaiorange}{\textbf{MEDIA}}%
    \else\textcolor{visaigreen}{\textbf{BASSA}}%
    \fi\fi
}
\newcommand{\visai}{\textsc{VisAi}}

% --- Nuova colonna per tabularx ---
\newcolumntype{L}{>{\raggedright\arraybackslash}X}
\newcolumntype{C}{>{\centering\arraybackslash}X}

% =============================================================================
\begin{document}

% --- Pagina di copertina ---
\begin{titlepage}
    \centering
    \vspace*{3cm}
    {\Huge\bfseries\textcolor{visaidark}{VisAi}\par}
    \vspace{0.5cm}
    {\LARGE\textcolor{visaiblue}{Analisi dei Requisiti Tecnici}\par}
    \vspace{1cm}
    {\large Sistema Anti-Taccheggio basato su\\Computer Vision ed Edge AI\par}
    \vspace{2cm}
    {\large\textcolor{visaigray}{Versione 1.0}\par}
    \vspace{0.5cm}
    {\large\textcolor{visaigray}{Luglio 2026}\par}
    \vspace{3cm}
    \begin{tabular}{ll}
        \textbf{Autore:} & Andrea Masiero \\
        \textbf{Progetto:} & VisAi — Computer Vision PoC \\
        \textbf{Classificazione:} & Confidenziale \\
    \end{tabular}
    \vfill
\end{titlepage}

% --- Indice ---
\tableofcontents
\newpage

% =============================================================================
\section{Introduzione}
% =============================================================================

\subsection{Scopo del Documento}

Il presente documento definisce i \textbf{requisiti tecnici} del progetto \visai{}, un sistema di prevenzione dei furti basato su Computer Vision ed Edge AI per il settore retail dell'abbigliamento. L'obiettivo è fornire una specifica completa e strutturata che guidi lo sviluppo del prodotto, dalla fase di Proof of Concept (PoC) fino al deployment in produzione.

\subsection{Visione del Progetto}

\visai{} mira a \textbf{sostituire completamente i sistemi antitaccheggio fisici} (placche EAS, antenne acustiche/magnetoacustiche e radiofrequenza) con una soluzione basata interamente sull'intelligenza artificiale. Il sistema è progettato per eliminare i falsi allarmi e i costi operativi legati alla fornitura, applicazione e rimozione delle placche fisiche.

\subsection{Target di Mercato}

Grandi catene di Retail, con focus iniziale su:
\begin{itemize}
    \item Fast Fashion (es.\ Zara, H\&M, Primark)
    \item Grandi magazzini (es.\ Rinascente, Coin)
    \item Catene sportswear (es.\ Decathlon, Nike Store)
\end{itemize}

\subsection{Glossario}

\begin{tabularx}{\textwidth}{lL}
    \toprule
    \textbf{Termine} & \textbf{Definizione} \\
    \midrule
    Edge AI & Esecuzione di modelli AI direttamente su hardware locale (on-premise), senza invio dati al cloud \\
    Action Recognition & Task di CV che riconosce e classifica azioni umane in un video \\
    Pose Estimation & Individuazione delle articolazioni chiave (keypoints) del corpo umano \\
    Object Counting & Conteggio automatico di oggetti in un flusso video \\
    Re-Identification & Riconoscimento della stessa entità attraverso viste o momenti temporali diversi \\
    EAS & Electronic Article Surveillance — sistemi antitaccheggio tradizionali \\
    FPS & Frames Per Second — frequenza di elaborazione dei fotogrammi \\
    ONNX & Open Neural Network Exchange — formato interoperabile per modelli AI \\
    TensorRT & Framework NVIDIA per ottimizzazione dell'inferenza su GPU \\
    \bottomrule
\end{tabularx}

% =============================================================================
\section{Architettura di Sistema}
% =============================================================================

\subsection{Panoramica Architetturale}

Il sistema adotta un'architettura \textbf{Edge AI pura}: ogni punto vendita è dotato di un server fisico proprietario che processa localmente tutti i flussi video. Nessun dato video viene trasferito al cloud.

\begin{center}
\begin{tikzpicture}[
    node distance=1.5cm and 2cm,
    box/.style={rectangle, draw=visaiblue, fill=visailight, rounded corners, 
                minimum width=3.2cm, minimum height=1cm, align=center,
                font=\small},
    server/.style={rectangle, draw=visaidark, fill=visaidark!10, rounded corners,
                   minimum width=3.5cm, minimum height=1.2cm, align=center,
                   font=\small\bfseries},
    alert/.style={rectangle, draw=visaired, fill=visaired!10, rounded corners,
                  minimum width=3cm, minimum height=1cm, align=center,
                  font=\small},
    arr/.style={-{Stealth[length=6pt]}, thick, color=visaigray}
]
    % Nodi
    \node[box] (cam1) {Telecamere\\Area Vendita};
    \node[box, below=of cam1] (cam2) {Telecamera\\Camerini};
    \node[server, right=2.5cm of $(cam1)!0.5!(cam2)$] (edge) {Server Edge\\Locale};
    \node[box, right=2.5cm of edge, yshift=1cm] (ar) {Action\\Recognition};
    \node[box, right=2.5cm of edge, yshift=-1cm] (oc) {Object\\Counting};
    \node[alert, right=2cm of $(ar)!0.5!(oc)$] (alert) {Sistema\\Alert};
    
    % Frecce
    \draw[arr] (cam1) -- (edge);
    \draw[arr] (cam2) -- (edge);
    \draw[arr] (edge) -- (ar);
    \draw[arr] (edge) -- (oc);
    \draw[arr] (ar) -- (alert);
    \draw[arr] (oc) -- (alert);
\end{tikzpicture}
\end{center}

\subsection{Principi Architetturali}

\begin{enumerate}
    \item \textbf{Privacy by Design}: nessun dato video esce dal punto vendita
    \item \textbf{Latenza zero}: processing locale per risposta in tempo reale
    \item \textbf{Resilienza}: funzionamento garantito anche offline (senza connessione internet)
    \item \textbf{Scalabilità orizzontale}: ogni negozio è un'unità indipendente
    \item \textbf{Hardware proprietario}: telecamere ottimizzate per il software VisAi
\end{enumerate}

% =============================================================================
\section{Requisiti Funzionali}
% =============================================================================

\subsection{RF-01: Rilevamento Anomalie Comportamentali (Action Recognition)}

\begin{tabularx}{\textwidth}{lL}
    \toprule
    \textbf{Campo} & \textbf{Specifica} \\
    \midrule
    \textbf{ID} & \reqid{RF-01} \\
    \textbf{Priorità} & \priority{H} \\
    \textbf{Descrizione} & Il sistema deve rilevare anomalie comportamentali e atteggiamenti sospetti finalizzati all'occultamento intenzionale di capi di abbigliamento. \\
    \textbf{Input} & Stream video in tempo reale dalle telecamere dell'area vendita \\
    \textbf{Output} & Alert con anomaly score, timestamp, ID telecamera e frame di riferimento \\
    \textbf{Vincolo critico} & Riduzione a zero dei falsi positivi: il sistema deve distinguere tra atti di furto e gesti normali (mettere le mani in tasca, sistemare la propria borsa, provare un capo) \\
    \bottomrule
\end{tabularx}

\subsubsection{Pipeline di Elaborazione}

\begin{enumerate}
    \item \textbf{Person Detection}: rilevamento di ogni persona nell'inquadratura
    \item \textbf{Pose Estimation}: stima dello scheletro (keypoints) per ogni persona
    \item \textbf{Skeleton Sequence Analysis}: analisi temporale della sequenza di movimenti
    \item \textbf{Action Classification}: classificazione dell'azione (normale / sospetta)
    \item \textbf{Anomaly Scoring}: assegnazione di un punteggio di anomalia
    \item \textbf{Alert Generation}: generazione dell'alert se lo score supera la soglia
\end{enumerate}

\subsubsection{Modelli Candidati}

\subsubsection{Sfide Tecniche}

\begin{itemize}
    \item \textbf{Zero falsi positivi}: distinguere tra furto e gesti quotidiani innocui
    \item \textbf{Occlusioni}: soggetti parzialmente nascosti da scaffali, manichini o altri clienti
    \item \textbf{Varietà di azioni}: l'occultamento può avvenire in modi molto diversi (in tasca, in borsa, sotto vestiti)
    \item \textbf{Inferenza real-time}: l'intera pipeline deve girare su Edge AI con latenza $< 500$~ms
    \item \textbf{Illuminazione variabile}: condizioni di luce diverse nelle varie zone del negozio
    \item \textbf{Multi-persona}: gestione simultanea di più soggetti nell'inquadratura
\end{itemize}

% ---

\subsection{RF-02: Conteggio Capi nei Camerini (Object Counting / Re-ID)}

\begin{tabularx}{\textwidth}{lL}
    \toprule
    \textbf{Campo} & \textbf{Specifica} \\
    \midrule
    \textbf{ID} & \reqid{RF-02} \\
    \textbf{Priorità} & \priority{H} \\
    \textbf{Descrizione} & Il sistema deve monitorare gli ingressi/uscite dei camerini, verificando la corrispondenza esatta tra il numero di capi portati all'interno e quelli portati fuori. \\
    \textbf{Input} & Stream video dalla telecamera dedicata all'ingresso dei camerini \\
    \textbf{Output} & Alert quando $\Delta = \text{capi\_in} - \text{capi\_out} > 0$ (capi mancanti) \\
    \textbf{Vincolo critico} & Non è ammesso filmare l'interno del camerino — il sistema opera esclusivamente sulle immagini di entrata e uscita \\
    \bottomrule
\end{tabularx}

\subsubsection{Pipeline di Elaborazione}

\begin{enumerate}
    \item \textbf{Person Detection + Tracking}: rilevamento e tracciamento della persona
    \item \textbf{Object Detection}: individuazione dei capi di abbigliamento trasportati (in mano, sul braccio, in borsa)
    \item \textbf{Count IN}: conteggio dei capi al momento dell'ingresso nel camerino
    \item \textbf{Count OUT}: conteggio dei capi al momento dell'uscita dal camerino
    \item \textbf{$\Delta$ Calculation}: calcolo della differenza $\Delta = \text{IN} - \text{OUT}$
    \item \textbf{Alert}: se $\Delta > 0$, generazione dell'alert (capo mancante)
\end{enumerate}

\subsubsection{Sfide Tecniche}

\begin{itemize}
    \item \textbf{Privacy del camerino}: nessuna immagine dell'interno — solo entrata/uscita
    \item \textbf{Occlusioni}: capi piegati, sovrapposti o nascosti in borse
    \item \textbf{Variabilità dei capi}: colori, dimensioni e textures estremamente variabili
    \item \textbf{Capi indossati}: il cliente può indossare un capo sotto i propri vestiti
    \item \textbf{Multi-persona}: più clienti che entrano/escono contemporaneamente
    \item \textbf{Capi restituiti allo staff}: il camerino può restituire capi non acquistati tramite un addetto
\end{itemize}

% ---

\subsection{RF-03: Sistema di Alert e Notifiche}

\begin{tabularx}{\textwidth}{lL}
    \toprule
    \textbf{Campo} & \textbf{Specifica} \\
    \midrule
    \textbf{ID} & \reqid{RF-03} \\
    \textbf{Priorità} & \priority{H} \\
    \textbf{Descrizione} & Il sistema deve generare alert in tempo reale e notificare il personale di sicurezza del punto vendita quando viene rilevata un'anomalia. \\
    \textbf{Input} & Anomaly score da RF-01 e/o $\Delta > 0$ da RF-02 \\
    \textbf{Output} & Notifica push, alert su dashboard locale, eventuale segnale acustico discreto \\
    \bottomrule
\end{tabularx}

\subsection{RF-04: Dashboard di Monitoraggio Locale}

\begin{tabularx}{\textwidth}{lL}
    \toprule
    \textbf{Campo} & \textbf{Specifica} \\
    \midrule
    \textbf{ID} & \reqid{RF-04} \\
    \textbf{Priorità} & \priority{M} \\
    \textbf{Descrizione} & Il sistema deve fornire una dashboard locale installata sul server Edge che mostri lo stato in tempo reale delle telecamere, gli alert attivi e lo storico degli eventi. \\
    \textbf{Input} & Dati da tutti i moduli AI e stream telecamere \\
    \textbf{Output} & Interfaccia web locale accessibile dalla rete del negozio \\
    \bottomrule
\end{tabularx}

\subsection{RF-05: Oscuramento Volti On-the-Edge}

\begin{tabularx}{\textwidth}{lL}
    \toprule
    \textbf{Campo} & \textbf{Specifica} \\
    \midrule
    \textbf{ID} & \reqid{RF-05} \\
    \textbf{Priorità} & \priority{H} \\
    \textbf{Descrizione} & Il sistema deve oscurare automaticamente i volti di tutte le persone riprese prima di qualsiasi registrazione o visualizzazione, garantendo la conformità all'EU AI Act. \\
    \textbf{Input} & Frame video grezzi \\
    \textbf{Output} & Frame con volti oscurati (blur/pixelation) prima di storage o display \\
    \textbf{Vincolo} & L'oscuramento deve avvenire on-the-edge, prima che il frame raggiunga qualsiasi storage \\
    \bottomrule
\end{tabularx}

% =============================================================================
\section{Requisiti Non Funzionali}
% =============================================================================

\subsection{RNF-01: Prestazioni e Latenza}

\begin{tabularx}{\textwidth}{lL}
    \toprule
    \textbf{Metrica} & \textbf{Requisito} \\
    \midrule
    Latenza end-to-end (evento $\rightarrow$ alert) & $< 2$ secondi \\
    Frame rate di elaborazione & $\geq 15$ FPS per telecamera \\
    Numero di telecamere simultanee & $\geq 8$ per server Edge \\
    Throughput totale & $\geq 120$ FPS aggregati su tutte le telecamere \\
    Tempo di avvio del sistema & $< 60$ secondi \\
    \bottomrule
\end{tabularx}

\subsection{RNF-02: Affidabilità e Disponibilità}

\begin{tabularx}{\textwidth}{lL}
    \toprule
    \textbf{Metrica} & \textbf{Requisito} \\
    \midrule
    Uptime & $\geq 99.5\%$ durante l'orario di apertura del negozio \\
    Recovery da crash & Restart automatico entro 30 secondi \\
    Funzionamento offline & Piena operatività anche senza connessione internet \\
    Gestione failure telecamera & Il sistema deve continuare a operare con le telecamere rimanenti \\
    \bottomrule
\end{tabularx}

\subsection{RNF-03: Accuratezza dei Modelli AI}

\begin{tabularx}{\textwidth}{lL}
    \toprule
    \textbf{Metrica} & \textbf{Requisito (target)} \\
    \midrule
    False Positive Rate (Action Recognition) & $< 1\%$ (obiettivo: $\sim 0\%$) \\
    Recall (Action Recognition) & $\geq 85\%$ \\
    Accuratezza conteggio capi (Object Counting) & $\geq 95\%$ \\
    Precision Person Detection & $\geq 98\%$ \\
    Accuratezza Face Blurring & $100\%$ (nessun volto non oscurato) \\
    \bottomrule
\end{tabularx}

\subsection{RNF-04: Sicurezza e Privacy}

\begin{tabularx}{\textwidth}{lL}
    \toprule
    \textbf{Requisito} & \textbf{Dettaglio} \\
    \midrule
    Nessun trasferimento video al cloud & Tutti i dati video restano sul server locale \\
    Crittografia a riposo & I dati salvati localmente devono essere cifrati (AES-256) \\
    Accesso autenticato & Dashboard e API protette da autenticazione (almeno username/password + HTTPS) \\
    Log di accesso & Audit trail di tutti gli accessi al sistema \\
    Data retention & Cancellazione automatica dei dati video dopo un periodo configurabile (default: 72h) \\
    \bottomrule
\end{tabularx}

\subsection{RNF-05: Scalabilità e Manutenzione}

\begin{tabularx}{\textwidth}{lL}
    \toprule
    \textbf{Requisito} & \textbf{Dettaglio} \\
    \midrule
    Aggiornamento remoto dei modelli AI & OTA (Over-The-Air) update dei pesi dei modelli senza downtime \\
    Aggiornamento software & Deploy remoto di nuove versioni del software applicativo \\
    Monitoraggio centralizzato & Heartbeat e metriche di salute inviate periodicamente a un backend centrale (solo metadati, mai video) \\
    Configurazione per negozio & Parametri (soglie, numero telecamere, layout) configurabili per ogni punto vendita \\
    \bottomrule
\end{tabularx}

% =============================================================================
\section{Requisiti Hardware}
% =============================================================================

\subsection{RH-01: Server Edge}

\begin{tabularx}{\textwidth}{lL}
    \toprule
    \textbf{Componente} & \textbf{Specifica minima consigliata} \\
    \midrule
    GPU & NVIDIA RTX 4090 o equivalente (24~GB VRAM) — in alternativa: NVIDIA L4 / T4 per data center \\
    CPU & Intel Core i7 / AMD Ryzen 7 (o superiore), $\geq 8$ core \\
    RAM & $\geq 32$~GB DDR5 \\
    Storage & $\geq 1$~TB NVMe SSD (per buffer video e log) \\
    Rete & Gigabit Ethernet (per connessione telecamere PoE e rete locale) \\
    Alimentazione & UPS integrato o esterno per protezione da blackout \\
    Form factor & Rack-mount o small form factor (installabile in armadio tecnico del negozio) \\
    Sistema operativo & Ubuntu 22.04 LTS Server (o superiore) \\
    \bottomrule
\end{tabularx}

\subsection{RH-02: Telecamere Proprietarie}

\begin{tabularx}{\textwidth}{lL}
    \toprule
    \textbf{Specifica} & \textbf{Requisito} \\
    \midrule
    Risoluzione & $\geq 1080$p (Full HD), consigliato 4K per area vendita \\
    Frame rate & $\geq 30$ FPS \\
    Lente & Grandangolare ($\geq 120°$ FoV) per area vendita; ottica fissa per camerini \\
    Visione notturna / IR & Sì, per condizioni di scarsa illuminazione \\
    Connettività & PoE (Power over Ethernet) — alimentazione e dati su un unico cavo \\
    Formato stream & RTSP / ONVIF compatibile \\
    Resistenza & IP65+ per telecamere esterne; IP20 per telecamere interne \\
    \bottomrule
\end{tabularx}

\subsection{RH-03: Configurazione Tipo per Punto Vendita}

\begin{tabularx}{\textwidth}{clL}
    \toprule
    \textbf{Zona} & \textbf{N. Telecamere} & \textbf{Funzione} \\
    \midrule
    Area vendita (piano principale) & $4\text{--}6$ & Action Recognition — copertura scaffali e zone ad alto rischio \\
    Ingresso camerini & $1\text{--}2$ & Object Counting — conteggio capi in/out \\
    Ingresso/uscita negozio & $1\text{--}2$ & Monitoraggio flussi e correlazione con alert \\
    \midrule
    \textbf{Totale per negozio medio} & $\mathbf{6\text{--}10}$ & \\
    \bottomrule
\end{tabularx}

% =============================================================================
\section{Requisiti Software e Stack Tecnologico}
% =============================================================================

\subsection{RS-01: Stack di Sviluppo AI}

\begin{tabularx}{\textwidth}{lL}
    \toprule
    \textbf{Componente} & \textbf{Tecnologia} \\
    \midrule
    Framework deep learning & PyTorch $\geq 2.x$ \\
    Inferenza ottimizzata & NVIDIA TensorRT / ONNX Runtime \\
    Object Detection & Ultralytics YOLOv8/v9 o RT-DETR \\
    Pose Estimation & HRNet / ViTPose (top-down) o OpenPose (bottom-up) \\
    Multi-Object Tracking & ByteTrack / BoT-SORT \\
    Action Recognition & ST-GCN / PoseC3D (skeleton-based) o SlowFast (3D CNN) \\
    Face Detection + Blurring & RetinaFace / SCRFD + Gaussian Blur \\
    \bottomrule
\end{tabularx}

\subsection{RS-02: Stack Applicativo}

\begin{tabularx}{\textwidth}{lL}
    \toprule
    \textbf{Componente} & \textbf{Tecnologia} \\
    \midrule
    Linguaggio principale & Python 3.11+ \\
    Video pipeline / decoding & FFmpeg, GStreamer, OpenCV \\
    API server locale & FastAPI \\
    Dashboard locale & React / Next.js (web-based, accessibile via browser) \\
    Database locale & SQLite o PostgreSQL (per eventi, alert, log) \\
    Containerizzazione & Docker + Docker Compose \\
    Orchestrazione & Systemd (per servizi) o Kubernetes K3s (per deploy più complessi) \\
    \bottomrule
\end{tabularx}

\subsection{RS-03: Ottimizzazione per Edge}

\begin{tabularx}{\textwidth}{lL}
    \toprule
    \textbf{Tecnica} & \textbf{Descrizione} \\
    \midrule
    Model Quantization & Riduzione precisione dei pesi: FP32 $\rightarrow$ FP16 $\rightarrow$ INT8 \\
    Pruning & Rimozione delle connessioni ridondanti dalla rete neurale \\
    Knowledge Distillation & Training di un modello compatto (student) da un modello più grande (teacher) \\
    TensorRT Compilation & Compilazione del modello ONNX in un engine TensorRT ottimizzato per la GPU specifica \\
    Batch Processing & Elaborazione di più frame simultaneamente per massimizzare il throughput GPU \\
    Model Caching & Pre-loading dei modelli in VRAM all'avvio per eliminare il cold-start \\
    \bottomrule
\end{tabularx}

% =============================================================================
\section{Requisiti di Compliance e Normativa}
% =============================================================================

\subsection{RC-01: GDPR (Regolamento UE 2016/679)}

\begin{itemize}
    \item \textbf{Minimizzazione dei dati}: il sistema processa solo i dati strettamente necessari
    \item \textbf{Nessun trasferimento}: i dati video non lasciano mai il server locale del punto vendita
    \item \textbf{Data retention limitata}: cancellazione automatica dopo il periodo configurato
    \item \textbf{Privacy by Design}: l'architettura Edge è stata scelta specificamente per la compliance
    \item \textbf{DPIA}: necessaria una Data Protection Impact Assessment prima del deployment
\end{itemize}

\subsection{RC-02: EU AI Act (Regolamento UE 2024/1689)}

\begin{itemize}
    \item \textbf{Classificazione del rischio}: il sistema potrebbe rientrare nei sistemi ad ``alto rischio'' (sorveglianza biometrica). L'oscuramento dei volti on-the-edge è la strategia per mitigare questo rischio
    \item \textbf{Oscuramento volti obbligatorio}: nessun dato biometrico viene estratto, memorizzato o trasmesso
    \item \textbf{Trasparenza}: obbligo di informare i clienti della presenza del sistema (cartellonistica)
    \item \textbf{Supervisione umana}: il sistema genera alert per il personale, non agisce autonomamente
\end{itemize}

\subsection{RC-03: Ispettorato del Lavoro e Sindacati}

\begin{itemize}
    \item \textbf{Art. 4 Statuto dei Lavoratori}: il sistema riprende aree frequentate anche da dipendenti — necessaria autorizzazione dell'Ispettorato del Lavoro o accordo sindacale
    \item \textbf{Finalità legittima}: la videosorveglianza deve avere finalità di sicurezza patrimoniale, non di controllo dei dipendenti
    \item \textbf{Informativa ai dipendenti}: obbligo di informare i lavoratori sull'esistenza e le finalità del sistema
    \item \textbf{Strategia VisAi}: fornire al cliente un ``pacchetto compliance'' pre-confezionato con template per le autorizzazioni
\end{itemize}

% =============================================================================
\section{Matrice di Tracciabilità dei Requisiti}
% =============================================================================

\begin{center}
\footnotesize
\begin{tabularx}{\textwidth}{llcL}
    \toprule
    \textbf{ID} & \textbf{Nome} & \textbf{Priorità} & \textbf{Dipendenze} \\
    \midrule
    \reqid{RF-01} & Action Recognition & \priority{H} & RH-01, RH-02, RS-01, RF-05 \\
    \reqid{RF-02} & Object Counting & \priority{H} & RH-01, RH-02, RS-01, RF-05 \\
    \reqid{RF-03} & Sistema Alert & \priority{H} & RF-01, RF-02, RS-02 \\
    \reqid{RF-04} & Dashboard & \priority{M} & RS-02, RH-01 \\
    \reqid{RF-05} & Face Blurring & \priority{H} & RS-01, RC-01, RC-02 \\
    \reqid{RNF-01} & Prestazioni & \priority{H} & RH-01, RS-03 \\
    \reqid{RNF-02} & Affidabilità & \priority{H} & RH-01, RS-02 \\
    \reqid{RNF-03} & Accuratezza AI & \priority{H} & RS-01, RF-01, RF-02 \\
    \reqid{RNF-04} & Sicurezza/Privacy & \priority{H} & RC-01, RC-02, RS-02 \\
    \reqid{RNF-05} & Manutenzione & \priority{M} & RS-02, RH-01 \\
    \reqid{RH-01} & Server Edge & \priority{H} & — \\
    \reqid{RH-02} & Telecamere & \priority{H} & RH-01 \\
    \reqid{RC-01} & GDPR & \priority{H} & RF-05, RNF-04 \\
    \reqid{RC-02} & EU AI Act & \priority{H} & RF-05 \\
    \reqid{RC-03} & Isp. Lavoro & \priority{H} & RC-01, RC-02 \\
    \bottomrule
\end{tabularx}
\end{center}

% =============================================================================
\section{Rischi Tecnici e Mitigazioni}
% =============================================================================

\begin{tabularx}{\textwidth}{clLL}
    \toprule
    \textbf{\#} & \textbf{Rischio} & \textbf{Impatto} & \textbf{Mitigazione} \\
    \midrule
    1 & Falsi positivi eccessivi & Perdita di fiducia del cliente, inutilizzabilità del sistema & Soglie conservative in produzione; training con dataset reali del negozio; periodo di ``shadow mode'' senza alert reali \\
    2 & Risorse GPU insufficienti & Frame drop, alert in ritardo & Ottimizzazione modelli (TensorRT, INT8); selezione di modelli leggeri; upgrade GPU \\
    3 & Occlusioni pesanti & Mancato rilevamento di azioni sospette & Multi-camera fusion; posizionamento strategico delle telecamere; modelli robusti alle occlusioni \\
    4 & Variabilità dell'illuminazione & Degradazione delle performance dei modelli & Data augmentation durante il training; telecamere con IR; preprocessing adattivo \\
    5 & Evoluzione normativa & Blocco del deployment per non conformità & Monitoraggio costante delle normative; design modulare per adattarsi rapidamente \\
    6 & Resistenza del personale & Rifiuto o boicottaggio del sistema & Comunicazione trasparente; coinvolgimento dei sindacati fin dalla fase pilota \\
    \bottomrule
\end{tabularx}

% =============================================================================
\section{Metriche di Successo (KPI)}
% =============================================================================

\begin{tabularx}{\textwidth}{lLl}
    \toprule
    \textbf{KPI} & \textbf{Descrizione} & \textbf{Target} \\
    \midrule
    False Positive Rate & Percentuale di alert generati erroneamente & $< 1\%$ \\
    Detection Rate & Percentuale di furti effettivi rilevati dal sistema & $\geq 85\%$ \\
    Latenza alert & Tempo tra l'evento e la ricezione dell'alert & $< 2$~s \\
    Accuratezza conteggio & Corrispondenza tra conteggio reale e stimato dei capi & $\geq 95\%$ \\
    Uptime & Disponibilità del sistema durante l'orario di apertura & $\geq 99.5\%$ \\
    ROI per negozio & Riduzione delle perdite per furto rispetto al costo del sistema & $> 3\times$ \\
    Tempo di deployment & Tempo per installare e configurare un nuovo punto vendita & $< 1$ giorno \\
    \bottomrule
\end{tabularx}

% =============================================================================
\section{Prossimi Passi}
% =============================================================================

\begin{enumerate}
    \item \textbf{PoC (Proof of Concept)}: implementare la pipeline Action Recognition con un singolo stream video su hardware di sviluppo, utilizzando un dataset pubblico di azioni anomale.
    \item \textbf{Dataset proprietario}: raccogliere un dataset di occultamento merce in ambiente controllato per il fine-tuning dei modelli.
    \item \textbf{Benchmark modelli}: confrontare le performance di SlowFast vs ST-GCN vs PoseC3D su hardware Edge target.
    \item \textbf{Prototipo Object Counting}: sviluppare la pipeline YOLO + ByteTrack + line crossing per il conteggio capi ai camerini.
    \item \textbf{Negozio pilota}: definire metriche, KPI e accordi per il primo test in un negozio reale.
\end{enumerate}

% =============================================================================
\vfill
\begin{center}
    \textcolor{visaigray}{\rule{0.5\textwidth}{0.4pt}}\\
    \vspace{0.5cm}
    \textcolor{visaigray}{\footnotesize Fine del documento — VisAi, Analisi dei Requisiti Tecnici v1.0}
\end{center}

\end{document}
